We measure CBA content similarity using four complementary metrics, each capturing a different dimension of agreement between clause vectors.
Let $\mathbf{x} \in \mathbb{R}^K$ denote an untreated firm's clause vector (the $K = 139$ clause-type counts in a given CBA period) and $\mathbf{y} \in \mathbb{R}^K$ the corresponding reference vector of connected treated firms.
\textbf{Cosine similarity}, $\cos(\mathbf{x}, \mathbf{y}) = \mathbf{x} \cdot \mathbf{y} \,/\, (\|\mathbf{x}\| \|\mathbf{y}\|)$, measures the angle between vectors and is invariant to the overall scale of clauses, so it captures agreement in \emph{composition} rather than volume.
\textbf{Bray-Curtis similarity}, $1 - \sum_k |y_k - x_k| / \sum_k (x_k + y_k)$, penalises absolute count differences relative to total clause mass; it is sensitive to changes in both the presence and the intensity of individual clause types.
\textbf{Total-variation similarity}, $1 - \tfrac{1}{2}\sum_k |\tilde{y}_k - \tilde{x}_k|$ where $\tilde{x}_k = x_k / \sum_k x_k$, converts each vector to a probability distribution over clause types and measures the overlap of those distributions; it equals one minus the total-variation distance and is therefore bounded in $[0,1]$.
\textbf{Ruzicka similarity}, $\sum_k \min(x_k, y_k) \,/\, \sum_k \max(x_k, y_k)$, is the count-weighted analogue of the Jaccard index: it measures the share of clause volume that is matched across the two vectors, with every unmatched unit on either side acting as a full penalty. It coincides with Jaccard when counts are binarised and shares the same numerator as Bray-Curtis but with a stricter (union-based) denominator, so $1 - $ Ruzicka is a proper metric.
All four measures equal one for identical vectors and zero for fully disjoint ones.
